\documentclass{article}
\usepackage{graphicx} % Required for inserting images
\usepackage{setspace}
\usepackage{amsmath}
\usepackage{ amssymb }
\usepackage{amsfonts} % for "\mathbb" macro
\usepackage{amsfonts}
\usepackage{tikz}
\usepackage{float}
\usepackage[a4paper, total={6in, 8in}]{geometry}

\title{Analysis and Design of Algorithms Lab\\Project Report}
\author{Ahmed Mohammed}
\date{May 2026}

\begin{document}

\parindent=0mm
\maketitle

\section{Problem Statement \& Application Context}
\subsection{Real-World Problem}
 
Competitive programming platforms such as Codeforces and LeetCode expose
users to thousands of problems but provide no personalized guidance on
what to solve next. As a result, users often waste limited study time on
problems that are too easy, too difficult, or unrelated to their current
weaknesses. This project proposes an algorithmic scheduling feature that
automatically generates an optimal, personalized study session given a
user's available time, skill level, and past performance.
 
\subsection{Formal Problem Definition}
 
\textbf{Input:}
\begin{itemize}
\item A problem set $P = \{p_1, p_2, \ldots, p_n\}$, where each
    problem $p_i$ carries a unique ID, a difficulty rating, an
        estimated solving time $t_i$ (minutes), and one or more topic
        tags.
  \item A user profile consisting of a skill level $s$ and a
        proficiency map from the set of topics to $[0, 1]$.
  \item A history of previous attempts, each recording whether the
        problem was solved, the time taken, and the number of wrong
        submissions.
  \item A session time budget $T$ (minutes).
\end{itemize}

\textbf{Output:} A subset $S \subseteq P$ such that
\[
  \sum_{p_i \in S} t_i \leq T
  \quad \text{and} \quad
  \sum_{p_i \in S} \mathrm{benefit}(p_i, \text{user})
  \text{ is maximized.}
\],

where $t_i$ is the time of problem $i$.

\textbf{Benefit function.} For a problem $p$ and user $u$, the benefit
is defined as:
\[
  \mathrm{benefit}(p, u)
  = \underbrace{\bar{w}(p,u)}_{\text{topic weakness}}
  \times \underbrace{e^{-0.5\,|d_p - s|}}_{\text{difficulty match}}
  \times \underbrace{\frac{1}{1+\text{attempts}(p,u)}}_{\text{recency penalty}}
  \times \underbrace{b(p,u)}_{\text{unsolved bonus}}
\]
where $\bar{w}(p,u)$ is the mean weakness across $p$'s topics
($1 - \pi(\text{topic})$), $d_p$ is the problem difficulty, $s$ is the
user's skill level, and $b(p,u) = 1.5$ if the user previously attempted
but failed $p$, $0.5$ if already solved, and $1.0$ otherwise.
 
% ─────────────────────────────────────────────
\section{Algorithmic Approaches}
% ─────────────────────────────────────────────
 
\subsection{Approach 1: Greedy}
 
\textbf{Design.} Each problem is assigned a priority score equal to its
benefit-to-time ratio $r_i = \mathrm{benefit}(p_i, \text{user}) / t_i$. Problems are
sorted in descending order of $r_i$; ties are broken by preferring
shorter problems. The scheduler then iterates through the sorted list
and greedily adds each problem to the session as long as it fits within
the remaining time budget.
 
\textbf{Justification.} The greedy ratio heuristic is the classical
approach for the fractional knapsack problem, where it yields the exact
optimum. The greedy algorithm actually performed very well in most test cases, compared to the second approach.
 
\textbf{Complexity.}
\begin{itemize}
  \item Time: $O(n \log n)$ dominated by sorting; the linear scan is
        $O(n)$.
  \item Space: $O(n)$ for the rated-problem list.
\end{itemize}
 
\textbf{Pros.}
\begin{itemize}
  \item Extremely fast.
  \item Memory space is negligible regardless of budget size.
\end{itemize}
 
\textbf{Cons.}
\begin{itemize}
  \item Does not guarantee the globally optimal selection for 0-1
        knapsack instances; it can leave large gaps of unused time.
\end{itemize}
 
\subsection{Approach 2: Dynamic Programming}
 
\textbf{Design.} Benefit values are real numbers; to apply integer DP
they are scaled by a constant $C = 10^6$ and truncated to integers.
A standard 0-1 knapsack DP table of size $(n+1) \times (T+1)$ is
filled, where $\mathrm{table}[i][t]$ stores the maximum scaled benefit
achievable using the first $i$ items with a time budget of $t$ minutes.
The selected problem set is recovered by backtracking through the table.
 
\textbf{Justification.} The 0-1 knapsack formulation models the problem
exactly: each problem is either included or excluded, estimated times
are positive integers, and benefit is additive. The DP solution is
guaranteed to be globally optimal (up to the scaling precision).
 
\textbf{Complexity.}
\begin{itemize}
  \item Time: $O(n \cdot T)$ where $T$ is the session budget in minutes.
  \item Space: $O(n \cdot T)$ for the DP table (stored as a flat
        \texttt{long long} array to avoid 2-D allocation overhead).
\end{itemize}
 
\textbf{Pros.}
\begin{itemize}
  \item Guarantees the optimal subset within the precision of the
        scaling constant.
\end{itemize}
 
\textbf{Cons.}
\begin{itemize}
  \item Memory and runtime grow with $T$; large budgets produce large tables.
  \item Scaling floating-point benefits to integers introduces a small error proportional to $1/C$.
  \item Significantly slower than greedy for large inputs.
\end{itemize}
 
% ─────────────────────────────────────────────
\section{Situational Evaluation}
% ─────────────────────────────────────────────
 
Both algorithms solve the same problem but make different trade-offs
between optimality and efficiency.
 
\textbf{When greedy is preferable.}
In interactive or real-time settings: for instance, a mobile app
that regenerates a session plan each time the user adjusts their time
budget, the $O(n \log n)$ greedy approach responds instantly even
for large problem sets. It is also preferable when the benefit values
of individual problems are close to each other, as the ratio ranking
tends to be near-optimal in that case. For typical competitive
programming sessions (budget of 60--180 minutes, problem times of
10--150 minutes), the greedy solution empirically achieves benefit
values within a few percent of the DP optimum.
 
\textbf{When knapsack DP is preferable.}
When the scheduler is run as a batch or background job (e.g.,
generating a weekly study plan overnight), the runtime cost of DP
is inconsequential, and the guaranteed optimality is worth it. DP
also outperforms greedy noticeably when problem times are large and
varied relative to the budget, creating scenarios where greedy's
inability to ``fill gaps'' leads to suboptimal selections.

% ─────────────────────────────────────────────
\section{Sample Test Cases}
% ─────────────────────────────────────────────

The following cases illustrate how the scheduler behaves under different
conditions. All cases use a time budget of $T = 124$ minutes.

\textbf{Case 1 : New user, no history.}
A user with skill level 1200 and no attempt history. Because no
recency penalty or unsolved bonus applies, benefit is driven purely by
topic weakness and difficulty match. Both algorithms mostly agree on the same
selection, since without history all benefits are smooth and the greedy
ratio ranking tends to be near-optimal.

\textbf{Case 2 : Experienced user with unsolved problems.}
A user with skill level 1200 and a history of 160 attempts, roughly
half of which are unsolved. Problems previously attempted but not solved
receive a $1.5\times$ unsolved bonus, making them high-priority
candidates. The knapsack DP exploits this by packing the session with
as many high-bonus problems as fit.

\textbf{Case 3 : Very short budget ($T = 20$ minutes).}
With a tight budget only the shortest high-benefit problems fit. The
greedy approach may leave a few minutes unused when no remaining
problem fits in the leftover time, while the DP always finds the
globally optimal fill. The benefit gap between the two algorithms is
most noticeable in this regime.

% ─────────────────────────────────────────────
\section{Performance Analysis}
% ─────────────────────────────────────────────

Table~\ref{tab:bench} and Figure~\ref{tab:graph} report measured runtimes and solution quality
across combinations of problem-set size and user history size, with a
fixed time budget of $T = 124$ minutes. Each runtime reading is the
average of five independent runs, reported in microseconds ($\mu$s).
The quality ratio column shows greedy total benefit as a percentage of
knapsack total benefit; a value of 100\% means both algorithms found
equally optimal solutions.

\begin{table}[h]
\centering
\caption{Runtime and solution quality: greedy vs.\ knapsack DP}
\label{tab:bench}
\begin{tabular}{cccccc}
\hline
\textbf{Problems} & \textbf{History} & \textbf{Knapsack} & \textbf{Greedy}
  & \textbf{Greedy/Knapsack} \\
 & \textbf{size} & \textbf{runtime ($\mu$s)} & \textbf{runtime ($\mu$s)}
  & \textbf{benefit (\%)} \\
\hline
250 & 55 & 695  & 621  & 99.65\% \\
250 & 45 & 641  & 628  & 99.65\% \\
250 & 30 & 1398 & 600  & 99.65\% \\
250 & 10 & 539  & 445  & 97.60\% \\
\hline
150 & 55 & 628  & 534  & 100.00\% \\
150 & 45 & 677  & 1174 & 100.00\% \\
150 & 30 & 707  & 439  & 100.00\% \\
150 & 10 & 479  & 419  & 100.00\% \\
\hline
50  & 55 & 229  & 179  & 100.00\% \\
50  & 45 & 219  & 168  & 99.79\% \\
50  & 30 & 208  & 150  & 99.79\% \\
50  & 10 & 146  & 90   & 99.79\% \\
\hline
\end{tabular}
\end{table}

\subsection*{Alignment with Theoretical Complexity}

\textbf{Greedy} has time complexity $O(n \log n)$, dominated by
sorting. This is visible in the data: as the problem count drops from
250 to 50, greedy runtime falls roughly proportionally. History size has negligible
effect on greedy runtime because the benefit computation is $O(h)$ per
problem (where $h$ is history size), but $h$ is small relative to the
sort cost.

\begin{figure}[H]
    \centering
    \includegraphics[width=0.75\linewidth]{image.png}
    \caption{Runtime: greedy vs. knapsack DP}
    \label{tab:graph}
\end{figure}

\textbf{Knapsack DP} has time complexity $O(n \cdot T)$, where $T$ is
the time budget. Since $T$ is fixed at 120 minutes across all runs, the
dominant variable is $n$. The measured runtimes decrease as $n$ shrinks
from 250 to 50, consistent with linear scaling in $n$. The weird
spike at (250 problems, history 30, 1398\,$\mu$s) is likely attributable
to CPU cache effects rather than algorithmic
behaviour, since the surrounding rows at the same problem count are
substantially lower.

\textbf{Solution quality.} The greedy algorithm matches the DP optimum
in the majority of runs (100\%), and never falls below 97.6\% of the
optimal benefit. This confirms that for this problem domain (moderate
budget, small estimated times, smooth benefit distribution) greedy is
a reliable approximation. The largest gaps appear at 250 problems with
history size 10, where fewer previously-attempted problems exist,
reducing the smoothing effect of the unsolved bonus and making the
benefit landscape less uniform.
 
% ─────────────────────────────────────────────
\section{Reflection}
% ─────────────────────────────────────────────

\subsection{AI Assistance}

AI tooling was used in two specific areas of this project.
First, the structure and components of the benefit function (topic weakness weighting, difficulty match score, recency penalty,
and unsolved bonus) were developed with AI suggestions.
Second, the input and output handling code (CSV parsing, file loading,
and writing routines across \texttt{csv\_utils}, \texttt{problem},
and \texttt{user}) was largely generated by AI and subsequently
corrected and integrated by hand to match the project's architecture.
All algorithmic logic (greedy selection and knapsack DP) was written
independently.

\subsection{Challenges}

\textbf{Modelling and file handling.} A significant portion of
development time was spent on system design: deciding how to represent
problems, user profiles, and attempt history; how to persist and load
them via CSV; and how the modules should interact. While file-handling
code was largely generated with AI assistance, considerable effort went
into reviewing, correcting, and integrating that code to match the
intended architecture. The core algorithmic work could have been
isolated and implemented much faster, but the goal was to build a
realistic, end-to-end prototype rather than a standalone algorithm demo.

\textbf{Handling floating-point values in DP.} The benefit function
produces real-valued scores, which are incompatible with the standard
integer knapsack DP. The chosen solution multiplies every benefit by a
large constant ($C = 10^6$) and truncates to an integer before filling
the DP table. This preserves the relative ordering of problems with
high precision while keeping the DP values bounded. The trade-off is a
small quantisation error: two problems whose true benefits differ by
less than $1/C$ may be treated as equal. This was deemed acceptable
given that the benefit function itself is already a heuristic
approximation.

\subsection{Limitations}

The prototype does not persist state between runs. User profiles and
problem data are static CSV files generated offline by dedicated
generator utilities (\texttt{gen\_problems}, \texttt{gen\_user},
\texttt{gen\_test}). In a real deployment the system would require a
database or API layer to track live solve events, update proficiency
scores incrementally, and synchronise with the platform's problem set.
Additionally, the benefit function is hand-crafted; a production system
might learn its weights from historical data.

\end{document}
